\documentclass{standalone}
\usepackage{circuitikz}
\usetikzlibrary{calc}
\definecolor{circuitnotemark}{HTML}{FE00FE}
\begin{document}
\begin{circuitikz}[american, line width=0.8pt]
\ctikzset{bipoles/length=1.2cm}
\ctikzset{grounds/scale=1.36}
\foreach \x in {0,2,4} {\foreach \y in {0,-2,-4,-6} {\fill[gray, opacity=0.35] (\x,\y) circle (1.1pt);}}
\node[gray, font=\scriptsize] at (0,0.84) {1};
\node[gray, font=\scriptsize] at (2,0.84) {2};
\node[gray, font=\scriptsize] at (4,0.84) {3};
\node[gray, font=\scriptsize, anchor=east] at (-0.84,0) {a};
\node[gray, font=\scriptsize, anchor=east] at (-0.84,-2) {b};
\node[gray, font=\scriptsize, anchor=east] at (-0.84,-4) {c};
\node[gray, font=\scriptsize, anchor=east] at (-0.84,-6) {d};
\coordinate (b1) at (0,-2);
\coordinate (d1) at (0,-6);
\coordinate (b2) at (2,-2);
\coordinate (b3) at (4,-2);
\coordinate (d3) at (4,-6);
\draw (b1) to[esource, l_=$V_{1}$, a^=$9\,\mathrm{V}$] (d1); % line 3
\draw (-0.1,-3.83) -- ++(0.2,0); % line 3
\draw (0,-3.93) -- ++(0,0.2); % line 3
\draw (-0.1,-4.17) -- ++(0.2,0); % line 3
\draw (b1) to[nos, l_=$S_{1}$] (b2); % line 4
\draw (b2) to[R, l_=$R_{1}$, a^=$10\,\mathrm{k}\Omega$] (b3); % line 5
\draw (b3) to[C, l_=$C_{1}$, a^=$100\,\mathrm{n}\mathrm{F}$] (d3); % line 6
\node[ground] at (d3) {}; % line 7
\draw (d1) -- (d3); % line 9
\node[circ] at (d3) {};
\path (8,-7.408) rectangle (14.181,0.593); % line 11
\node[anchor=west, inner sep=0, circuitnotemark, font=\large] at (8,0) {X}; % line 11
\node[anchor=west, inner sep=0, circuitnotemark, font=\large] at (8,-0.487) {X}; % line 11
\node[anchor=west, inner sep=0, circuitnotemark, font=\large] at (8,-0.974) {X}; % line 11
\node[anchor=west, inner sep=0, circuitnotemark, font=\large] at (8,-1.46) {X}; % line 11
\node[anchor=west, inner sep=0, circuitnotemark, font=\large] at (8,-1.947) {X}; % line 11
\node[anchor=west, inner sep=0, circuitnotemark, font=\large] at (8,-2.434) {X}; % line 11
\node[anchor=west, inner sep=0, circuitnotemark, font=\large] at (8,-2.921) {X}; % line 11
\node[anchor=west, inner sep=0, circuitnotemark, font=\large] at (8,-3.408) {X}; % line 11
\node[anchor=west, inner sep=0, circuitnotemark, font=\large] at (8,-3.895) {X}; % line 11
\node[anchor=west, inner sep=0, circuitnotemark, font=\large] at (8,-4.381) {X}; % line 11
\node[anchor=west, inner sep=0, circuitnotemark, font=\large] at (8,-4.868) {X}; % line 11
\node[anchor=west, inner sep=0, circuitnotemark, font=\large] at (8,-5.355) {X}; % line 11
\node[anchor=west, inner sep=0, circuitnotemark, font=\large] at (8,-5.842) {X}; % line 11
\node[anchor=west, inner sep=0, circuitnotemark, font=\large] at (8,-6.329) {X}; % line 11
\node[anchor=west, inner sep=0, circuitnotemark, font=\large] at (8,-6.816) {X}; % line 11
\node[anchor=south west, inner sep=2pt, circuitnotemark, font=\LARGE\bfseries] (circuittitle) at (current bounding box.north west) {X};
\path (circuittitle.south west) rectangle ++(6.858,0.853);
\end{circuitikz}
\end{document}
